\chapter{Modeling and Design of the Proposed Solution}

\section{Introduction}
Introduce the proposed application and define its official name or brand.

\section{Requirements Analysis}
\begin{itemize}
    \item Functional requirements
    \item Non-functional requirements
    \item Constraints and assumptions
\end{itemize}

\section{Global Architecture}
Describe the high-level architecture and module decomposition (frontend, backend, data, IAM, etc.).

\begin{figure}[H]
    \centering
    % \includegraphics[width=0.9\textwidth]{figures/architecture-overview.png}
    \caption{High-level architecture of the proposed solution}
    \label{fig:architecture-overview}
\end{figure}

\section{Modeling}
Include relevant models and diagrams:
\begin{itemize}
    \item use case diagrams,
    \item class or domain model,
    \item sequence or activity diagrams,
    \item database model if applicable.
\end{itemize}

\section{Design Decisions}
Explain key decisions and trade-offs:
\begin{itemize}
    \item chosen frameworks and patterns,
    \item security and authorization choices,
    \item scalability and maintainability considerations.
\end{itemize}

\section{Chapter Conclusion}
Summarize how the design responds to identified requirements and prepares implementation.

// my own words that claude will use later to generate the actual content of the chapter.


./path means the root of the opened folder (rapport-pfe)

#introduction
Canny TMS, a custom-built Ticket Management System designed to meet the specific needs of Canny Vision, was my responsability
to design and develop.
#requirements analysis
## functional requirements
Canny vision already has its own internal identity access managment software that handles authenticatoin and authorizatoin
that's based on oauth2.1 with pkce (dear claude read ./middleware/iam module for more details and to understand how the iam works in detail)
so the app that I developed had to be compatible with the existing IAM and use it for all authentication and authorization needs.
Canny vision also don't use any frameworks or libraries for the frontend, so the app had to be developed using vanilla js and no external dependencies.
that's why they have their own developed mvp engine (youc can read ./frontend/scripts/mvp.js to understand how it works) that I had to use for the frontend development.
the app also had to be developed using a zero trust architecture, which means that all security measures had to be implemented internally and no data could be stored on external servers.
## non functional requirements
the ui/ux design had to be simple and intuitive, and the app had to be responsive and accessible on different devices.
the app also had to be scalable and maintainable, as it would be used by multiple teams
and would need to evolve over time to meet changing requirements.
#global architecture
I will pass a figure here later that shows the high level architecture of the system
 read it yourself and wrtie what you see fit
#modeling
same thing too for global architecute
#design decisions

bootstrap vs tailwindcss
I started with bootsrap , which is the company choice (imported from cdn ; doesn't expose secuirty issues)
it's good because blah blah but it limits the design and makes it look like every other bootstrap app, and it also has a lot of unused css that affects performance.
I convinced Mr Kaaniche to switch to tailwindcss, which is a utility-first css framework that allows for more flexibility and better performance.
In normal methods (not so secure app) it should be installed via npm and used as a build tool, but since we can't use any external dependencies, I used it with standalone cli where I isntalled the engine locally
read : https://tailwindcss.com/blog/standalone-cli  and put it in the report in its place(idk the actual page where we should put from wehre we got our information)


abac: 
nist standard : I spent time reading and analyzing the nist standard for access control : here's the link https://nvlpubs.nist.gov/nistpubs/specialpublications/nist.sp.800-162.pdf
integrate it too in the report
to actually implenet it I used authZforce whihc is an open source xacml engine where I donwloade the source code and inspired from to implmeent my own

[I will put here the figure of nist compliant abac]
and saying how I didn't implement PAP because it's too next leven and maybe later 
